\documentclass[12pt, a4paper]{article}
\usepackage[T1]{fontenc}
\usepackage{lmodern}
\usepackage{amsmath, amsthm, amssymb, mathtools}
\usepackage[margin=1in]{geometry}
\usepackage{xr-hyper}
\usepackage[colorlinks=true,linkcolor=blue,citecolor=blue]{hyperref}
\usepackage{cleveref}
\usepackage{graphicx, subcaption}
\usepackage{natbib}
\usepackage{booktabs}
\usepackage{enumerate}

% Code references: scaled monospace with line-breaking
\newcommand{\code}[1]{\texttt{\small #1}}

\newtheorem{theorem}{Theorem}[section]
\newtheorem{lemma}[theorem]{Lemma}
\newtheorem{proposition}[theorem]{Proposition}
\newtheorem{corollary}[theorem]{Corollary}
\newtheorem{definition}{Definition}[section]
\newtheorem{remark}{Remark}[section]
\newtheorem{conjecture}[theorem]{Conjecture}
\newtheorem{derivation}[theorem]{Derivation}
\newtheorem{observation}[theorem]{Observation}
\numberwithin{equation}{section}

\newcommand{\dd}{\mathrm{d}}
\newcommand{\seriestitle}{Why Rotating Fluids Make Polygons}


\title{Why Rotating Fluids Make Polygons:\\[6pt]
\large A Companion for the Curious Reader}
\author{Gordon Speagle}
\date{}

\begin{document}
\maketitle

\section*{Prologue: The Polygon Number}

\subsection*{The hexagon that started this}

Saturn has a hexagon at its north pole. Not a rough hexagonal shape, not
a six-sided approximation---a persistent, symmetric, six-sided polygon
of clouds, each side longer than the diameter of Earth. The Voyager
spacecraft photographed it in 1981. The Cassini orbiter confirmed it
was still there thirty years later, rotating in lockstep with the
planet's interior. Jupiter has an octagon at its south pole: eight
cyclones arranged in a regular ring, each the size of the continental
United States, mapped by the Juno mission starting in 2017.

These are not weather in the way we think of weather. A hurricane on
Earth forms, wanders, and dies. These polygons persist for decades,
possibly centuries. They do not drift. They do not deform. They sit at
the poles and spin.

Why polygons? Why \emph{these} polygons---six for Saturn, eight for
Jupiter? I first encountered the Voyager photos years ago and did what
most people do: marveled, read a few explanations, and filed it under
``Nature Gonna Geometry.'' Crystals, snowflakes, beehives, the
nautilus shell---the physical world converges on geometric forms for
all sorts of reasons. I assumed the hexagon was another entry on that
list.

Years later I was watching an astronomy video with my kids and the
polar polygons came up again. My children were astonished. I was
curious. I dove back into the literature, ordered a copy of
\emph{Geophysical Fluid Dynamics} by Joseph Pedlosky, and started
writing code.

What I found was a formula---published in 1931 by a mathematician
named T.~H.~Havelock---that counts how many spinning objects can
sit in a stable ring. The formula is not complicated. It has two
parts: one from geometry, one from the surface the objects sit on.
The geometry part is universal; it does not care whether the objects
are storms on a gas giant, quantum vortices in superfluid helium, or
anything else that spins logarithmically. The surface part changes
from problem to problem.

On a flat surface, the formula says the largest stable ring has seven
members. Seven. Not six, not eight---seven. Saturn's hexagon has six
sides; Jupiter's octagon has eight. Both are one step away from seven,
on opposite sides.

I became determined to pull the connecting string until the knot
unraveled. The string led from fluid dynamics to algebraic number
theory to general relativity to particle physics to cosmology. I
followed it for years. The result is a six-paper technical series.
This companion document is my attempt to explain what that series
says, and why, without requiring any of the mathematics that makes
it work.


\subsection*{The one idea}

Here is the central claim, stated as plainly as I can manage.

When identical spinning objects arrange themselves in a ring on a
surface, there is a number---call it $N$---that counts how many
objects are in the ring. The ring's stability depends on a single
quantity that splits into two parts. The first part comes from the
surface: flat, curved like a sphere, curved like a saddle. The
second part comes from pure geometry: how many objects there are and
how they interact pairwise.

On a flat surface, the two parts balance exactly at $N = 7$. A ring
of seven is marginally stable---it neither grows nor shrinks when
disturbed. Below seven, the ring snaps back. Above seven, the ring
flies apart.

This number seven is not adjustable. It is not a parameter I chose or
a threshold I fit to data. It falls out of the formula the way four
falls out of the corners of a square. It is a property of the
two-dimensional logarithmic interaction---the mathematical law that
governs how point-like spinning objects push each other around on a
surface.

The technical papers argue that this same threshold determines far
more than atmospheric polygons. The number seven, and the algebraic
structure surrounding it, fixes the spin of the graviton (the
hypothetical particle that carries gravity); forces the gauge group
of particle physics to be $\mathrm{SU}(3) \times \mathrm{SU}(2)
\times \mathrm{U}(1)$ (the symmetry behind the strong, weak, and
electromagnetic forces); and constrains the cosmological constant
(the energy density of empty space). These are not separate
coincidences bolted together. They are consequences of one formula
applied to one threshold.

Whether the argument is correct is for the scientific community to
assess. My goal here is to make it accessible enough that you can
follow the logic, identify where the assumptions are, and form your
own judgment.


\subsection*{The scorecard}

The theory takes in one integer ($N = 7$), one scale (the Planck
mass, which sets the units of gravity), and three identifications
that connect the mathematical formalism to physics. From these
inputs it produces numerical predictions. The table below summarizes
the main ones.

\begin{center}
\small
\renewcommand{\arraystretch}{1.25}
\begin{tabular}{p{5.2cm} p{3.5cm} p{2cm} l}
\toprule
\textbf{What the theory predicts} & \textbf{What we measure} & \textbf{How close} & \textbf{Status} \\
\midrule
The gauge group of particle physics is $\mathrm{SU}(3) \times \mathrm{SU}(2) \times \mathrm{U}(1)$
  & Confirmed experimentally & Exact match & Theorem \\[4pt]
Three generations of matter particles
  & Three observed & Exact match & Derived \\[4pt]
The CKM phase (a parameter governing matter--antimatter asymmetry) is $70.2^\circ$
  & $73.5^\circ \pm 5.1^\circ$ & $0.4$ standard deviations & Derived \\[4pt]
The Weinberg angle (which sets the ratio of electromagnetic to weak force strength) gives $\sin^2\theta_W = 3/11$
  & Measured value differs by $3.9\%$ after quantum corrections & $3.9\%$ & Derived \\[4pt]
A quark mixing parameter $|V_{cb}| = 0.044$
  & Measured $\approx 0.042$ & $6\%$ & Derived \\[4pt]
The strange-to-bottom quark mass ratio is $0.023$
  & Measured $\approx 0.024$ & $4\%$ & Derived \\[4pt]
The Hubble constant (expansion rate of the universe) is $67.4$ km/s/Mpc
  & $67.4 \pm 0.5$ km/s/Mpc & $0.002\%$ & Derived \\[4pt]
Dark matter makes up $26.9\%$ of the cosmic energy budget
  & $26.8\% \pm 0.5\%$ & $0.4$ standard deviations & Conjectured \\[4pt]
Ordinary matter makes up $5.0\%$
  & $4.9\% \pm 0.1\%$ & $0.6$ standard deviations & Derived \\[4pt]
The Higgs boson mass is $125 \pm 3$ GeV
  & $125.25 \pm 0.17$ GeV & $2\%$ & Derived \\[4pt]
The strong coupling constant $\alpha_s = 0.115$--$0.122$
  & $0.1180 \pm 0.0009$ & Within range & Derived \\[4pt]
The string tension ratio $\sigma/\Lambda^2 = 5.97$
  & Lattice QCD: $6.25 \pm 0.5$ & $5\%$ & Derived \\[4pt]
The proton is exactly stable (never decays)
  & No decay observed & Consistent & Derived \\[4pt]
The QCD vacuum angle $\theta_{\mathrm{QCD}} = 0$
  & $< 10^{-10}$ & Consistent & Derived \\[4pt]
Neutrinos follow the normal mass hierarchy
  & Pending (JUNO experiment) & Future test & Derived \\[4pt]
Dark energy equation of state $w = -1$
  & $-1.0 \pm 0.1$ (DESI) & Consistent & Derived \\
\bottomrule
\end{tabular}
\end{center}

\noindent
A few things to notice. First, the ``How close'' column. Most
predictions match observations within a few percent or better. Second,
the ``Status'' column distinguishes between results that are
mathematically proved, results that are derived through a chain of
identifications, and results that remain conjectured. Third, several
predictions are not yet testable---they await experiments currently
under construction. Those are genuine predictions in the scientific
sense: stated before the measurement, falsifiable by the outcome.


\subsection*{The honest disclaimer}

I am not a physicist by training. My background is philosophy with a
mathematical bent; professionally, I am a software developer working
in healthcare. A hobbyist's obsession with number theory is an apt
description of my post-educational years. I am unaffiliated with any
university or research institution.

This matters because the claims in the technical series are
extraordinary, and extraordinary claims face a higher bar. I want to
be transparent about both the strengths and the limitations.

The results fall into three tiers of confidence:

\textbf{Proved.} A small number of results are mathematical theorems.
If you accept the starting identification---that a point vortex on a
two-dimensional surface and a gravitational source share the same
Green's function---then these results follow by logical necessity. The
gauge group of particle physics being
$\mathrm{SU}(3) \times \mathrm{SU}(2) \times \mathrm{U}(1)$ is in
this category. The proof uses finite group theory (specifically, the
symmetry group PSL(2,7) acting on the Klein quartic surface), and the
steps are checkable by anyone with an undergraduate algebra textbook.

\textbf{Derived.} Most results are in this tier. They follow from a
chain of identifications---each step uses standard physics or
standard mathematics---with controlled approximations whose errors
are estimated. The Weinberg angle, the Hubble constant, the quark
mixing parameters, and the Higgs mass are all derived. These results
could be wrong if one of the identifications in the chain fails, but
the chain is explicit: every link is written down, and each can be
tested independently.

\textbf{Conjectured.} A handful of results are numerically verified
but rely on at least one step that is physically motivated rather than
mathematically proved. The dark matter fraction is in this category.
The number matches, but the derivation has a gap that I have flagged
honestly.

Institutional research correctly forces disciplinary depth; the cost
is that cross-domain connections can go unnoticed. I am unaffiliated,
which meant I could follow the mathematics wherever it led, without
the professional incentive to stay within a single field. Whether that
freedom produced something worth the lack of expert supervision is for
the reader to assess.


\subsection*{The map}

The technical series has six papers. Here is what each one does.

\textbf{Paper I: The Mathematics.} This paper establishes the
foundation. The logarithmic interaction---the same law that governs
how spinning fluid parcels push each other---is the unique
two-dimensional interaction for which polygon stability is
scale-independent. The Havelock formula (1931) gives the stability
eigenvalues. On a flat surface the critical threshold is $N = 7$.
On curved surfaces the threshold shifts, and the shift values live
in algebraic number fields: the threshold for the transition from
seven to eight stable vortices is $8 - 3\sqrt{7}$, a unit in the
ring of integers of $\mathbb{Q}(\sqrt{7})$. This is where the
algebraic number theory enters, and it is not decoration---these
algebraic structures determine the physical predictions downstream.

\textbf{Paper II: The Physics.} This paper connects the mathematics
to actual fluid dynamics. Point vortices have no inertia; they are
pushed sideways by their neighbors. A central vortex can stabilize
rings that would otherwise be unstable. The paper derives the
critical central-vortex strength for each $N$, resolves the $N = 7$
marginal case with a quartic (fourth-order) calculation, and shows
that the stability results survive two realistic corrections: the
finite size of real vortices (the ``blob'' correction, which is
independent of the vortex profile shape) and the Rossby-wave
stationarity constraint that pins Saturn's hexagon at $N = 6$. The
final sections identify a structural isomorphism between point
vortices and gravitational sources: both produce conical
singularities with the same Green's function.

\textbf{Paper III: Gravity.} The marginal mode at $N = 7$ carries
spin 2---the same spin as the graviton, the hypothetical quantum of
gravity. This paper derives Einstein's equation (the field equation
of general relativity) from the statistical mechanics of vortex
configurations, using Jensen's inequality and the Onsager
large-energy limit. It constructs the Wheeler--DeWitt equation (the
quantum gravity wave equation) from the breathing mode of the
polygon, with the central charge $c = 12\,b(N)$ determined by the
Fisher--Rao metric on the space of vortex configurations. Black hole
thermodynamics, holography, and the emergence of quantum mechanics
from statistical mechanics are all addressed.

\textbf{Paper IV: Particle Physics.} The polygon hierarchy forces
the internal geometry to be an eleven-dimensional manifold (four
spacetime dimensions plus a seven-dimensional compact space). The
compact space is rigid: Thurston's geometrization theorem eliminates
all possibilities except one. From that geometry, the Standard Model
gauge group $\mathrm{SU}(3) \times \mathrm{SU}(2) \times
\mathrm{U}(1)$ is derived, not assumed. The Weinberg angle comes
out as $\sin^2\theta_W = 3/11$, a ratio of Casimir invariants. The
CKM phase, quark mass ratios, the Higgs mass, the strong coupling
constant, and the string tension of quark confinement are all
computed from the same central charge.

\textbf{Paper V: Cosmology.} The cosmological constant---the energy
density of empty space---is determined by the instanton action of the
$N = 11$ Born--Oppenheimer tunneling amplitude. The hierarchy between
the Planck scale and the electroweak scale (a factor of $10^{16}$)
is not put in by hand; it emerges from the tunneling exponent. The
cosmic energy budget (dark energy, dark matter, ordinary matter)
decomposes into tree-level, one-loop, and tunneling contributions.
The Hubble constant, the neutrino mass hierarchy, and the initial
conditions of inflation are derived.

\textbf{Paper VI: The Scorecard.} This paper collects every
prediction, classifies each by confidence tier, quantifies the
cumulative uncertainty, and lists six binary pass/fail experimental
tests---results from JUNO, LEGEND, Hyper-Kamiokande, LISA, and
DESI---that can confirm or falsify specific claims. It also
catalogues what remains open: the dark matter production mechanism,
the baryon asymmetry factor, and the light quark masses.


\newpage


\part{The Mathematics of Polygons}
% [Mirrors Paper I --- to be written]


\part{The Physics of Vortex Rings}
% [Mirrors Paper II --- to be written]


\part{How Gravity Emerges}
% [Mirrors Paper III --- to be written]


\part{The Particles Inside the Polygon}
% [Mirrors Paper IV --- to be written]


\part{The Numbers of the Universe}
% [Mirrors Paper V --- to be written]


\part{The Scorecard}
% [Mirrors Paper VI --- to be written]


\section*{Epilogue: One Formula}
% [To be written]


\end{document}
